\chapter{Dowody formalne i dedukcja naturalna}\label{ch:natural-deduction}

\section{Od dowodu pisanego słowami do obiektu matematycznego}
W rozdziale~\ref{ch:quantifiers} poznaliśmy typowe techniki dowodzenia,
a w rozdziałach~\ref{ch:syntax-fol} i~\ref{ch:semantics-fol}
oddzieliliśmy formuły od ich interpretacji. Teraz uczynimy
samą procedurę dowodzenia przedmiotem definicji.
Potrzebujemy skończonego obiektu, którego każdy krok da się sprawdzić
bez zgadywania, co autor „miał na myśli”. Taką rolę spełnia
\emph{formalna dedukcja}: drzewo formuł budowane z założeń
za pomocą jawnie wypisanych reguł wnioskowania.

Wybierzemy \emph{dedukcję naturalną}, ponieważ jej reguły
odpowiadają znanym metodom argumentowania. Inne systemy,
na przykład systemy Hilberta i rachunek sekwentów, omówimy
porównawczo w dalszej części.

\section{Jak czytać zapis reguły wnioskowania?}\label{sec:nd-how-to-read}
Zanim wypiszemy formalne reguły, nauczymy się czytać ich zapis.
To ważne: pozioma kreska, którą za chwilę zobaczymy, \emph{nie jest kreską ułamkową}
i nie oznacza dzielenia liczb. Oznacza przejście od podanych zdań do nowego zdania.

\begin{definicja}[Schemat reguły wnioskowania]\label{def:nd-rule-scheme}
\emph{Reguła wnioskowania} jest umową określającą, kiedy z już uzyskanych
formuł wolno zapisać kolejną formułę. Jej ogólny schemat ma postać
\[
 \frac{\text{przesłanka 1}\quad\cdots\quad\text{przesłanka }n}
      {\text{wniosek}}\;\text{(nazwa reguły)}.
\]
To, co zapisano \emph{nad kreską}, jest wymagane przed zastosowaniem reguły;
to, co zapisano \emph{pod kreską}, jest otrzymanym wnioskiem.
Przez \emph{przesłankę} rozumiemy tutaj formułę wcześniej uzyskaną w dowodzie
lub przyjętą jako założenie. Reguła może mieć jedną przesłankę, kilka przesłanek,
a czasem żadnej. Schemat opisuje wszystkie dopuszczalne użycia jednej reguły,
nie tylko jeden konkretny dowód.
\end{definicja}

\begin{przyklad}[Dwie przesłanki nad kreską]
Jeśli umiemy uzasadnić zdanie $A$ oraz zdanie $B$, to mamy prawo
wywnioskować ich koniunkcję $A\land B$. Zapis
\[
\frac{A\qquad B}{A\land B}\;(\land I)
\]
czytamy słownie: „\emph{z A oraz B wyprowadzamy A i B}”.
Wcześniej moglibyśmy np. ustalić, że liczba $6$ jest parzysta i że jest
dodatnia. Wtedy z tych dwóch ustaleń otrzymujemy zdanie
„liczba $6$ jest parzysta i dodatnia”. Nie sprawdzamy w tym kroku
na nowo obu faktów: korzystamy z tego, że zostały już uzasadnione.
Puste miejsce między $A$ i $B$ nad kreską oznacza tylko odstęp
w drzewie: \emph{obie} przesłanki są wymagane, to nie jest symbol
alternatywy. Gdybyśmy mieli tylko jedną przesłankę, tej reguły
zastosować nie wolno.
\end{przyklad}

\begin{uwaga}[Co oznaczają litery i skróty przy regułach?]
Litery $A,B,P,Q$ lub greckie $\varphi,\psi,\theta$ występujące w schematach
są \emph{metazmiennymi}: symbolizują dowolne formuły.
Nie należy ich utożsamiać ze zmiennymi $x,y$ występującymi wewnątrz formuł.
W nazwie $(\land I)$ litera $I$ pochodzi od angielskiego
\emph{introduction} („wprowadzenie” spójnika do wniosku),
a $E$ w $(\land E)$ od \emph{elimination} („eliminacja” spójnika).
Indeks $1$ lub $2$ przy $E$ mówi, którą część koniunkcji wybieramy:
$\land E_1$ daje pierwszy człon, $\land E_2$ drugi.
Analogicznie $\lor I_1$ i $\lor I_2$ wprowadzają alternatywę,
odwołując się odpowiednio do jej pierwszego lub drugiego członu.
\end{uwaga}

\begin{przyklad}[Reguła z jedną przesłanką]
Jeśli mamy już dowód formuły $A\land B$, to zapis
\[
\frac{A\land B}{A}\;(\land E_1)
\]
pozwala nam uzyskać $A$. To nie jest równoważność obu formuł:
mając $A$, nie możemy tą samą regułą uzyskać $A\land B$,
bo nie wiemy nic o $B$. Kreskę czytamy od góry do dołu,
a nie dowolnie w obu kierunkach. Na przykład z „$6$ jest liczbą
parzystą i dodatnią” można wywnioskować „$6$ jest liczbą parzystą”,
ale z samej parzystości rozważanej liczby nie wynika jej
dodatniość: na przykład $-6$ jest parzyste i ujemne.
Odwrotne przejście byłoby nieuprawnione.
\end{przyklad}

\begin{uwaga}[Drzewo dowodu]
Kilka reguł można zestawić w \emph{drzewo dowodu}.
Liście drzewa są u góry, a ostateczny wniosek na dole.
Kiedy jedna reguła kończy się na formule $A$, tę formułę można
wykorzystać niżej jako przesłankę innej reguły.
Symbol $\vdots$ (pionowe kropki) oznacza w schematach
pewien \emph{ciąg dalszych kroków}, które na razie pomijamy na rysunku;
nie jest osobną regułą ani formułą.
Nawiasy kwadratowe $[A]^i$ oznaczają założenie tymczasowe
o \emph{etykiecie} $i$. Nie orzekają, że $A$ jest prawdziwe.
Reguła z tym samym indeksem na dole może \emph{rozładować}
wybrane wystąpienia założenia, czyli sprawić, że końcowy wniosek
nie będzie już od niego zależał jako od otwartego założenia.
Wyjaśnimy to bezpośrednio na przykładach poniżej.
\end{uwaga}


\section{Co kryją trzy pionowe kropki?}\label{sec:nd-vertical-dots}
Najbardziej tajemniczym elementem drzewa dowodu może być symbol
$\vdots$. Przypomina trzy pionowo ustawione kropki i jest
\emph{skrótem graficznym}, a nie nowym spójnikiem, kwantyfikatorem,
regułą wnioskowania ani twierdzeniem. Informuje, że pomiędzy
wypisanymi formułami znajduje się fragment formalnego dowodu,
którego kolejnych kroków nie pokazano.

\begin{definicja}[Znaczenie symbolu $\vdots$ w drzewie dowodu]
Jeżeli zapisujemy fragment
\[
\begin{array}{c}
\mathcal D\\[-2pt]
\vdots\\[-2pt]
\varphi
\end{array}
\]
to symbol $\mathcal D$ oznacza wcześniej rozpoczętą dedukcję,
a $\vdots$ zastępuje skończone drzewo poprawnych zastosowań
reguł prowadzące do formuły $\varphi$.
Skrót jest dopuszczalny tylko wtedy, gdy takie drzewo
rzeczywiście istnieje, a jego założenia i warunki użycia
reguł są zgodne z otaczającą częścią dowodu.
Nie istnieje reguła pozwalająca z dowolnej formuły $\psi$
przejść za pomocą samych kropek do dowolnej formuły $\varphi$.
\end{definicja}

\begin{przyklad}[Od skrótu do wszystkich kroków]
Weźmy formułę $A\land B$. Jeżeli zamierzamy wyprowadzić
$B\land A$, możemy skrótowo napisać
\[
\begin{array}{c}
[A\land B]^1\\
\vdots\\
B\land A
\end{array}
\]
Jednak czytelnik ma prawo zapytać, \emph{jakimi regułami}
doszliśmy do ostatniej formuły. Rozwinięty fragment wygląda tak:
\[
\frac{
  \displaystyle\frac{[A\land B]^1}{B}\;(\land E_2)
  \qquad
  \displaystyle\frac{[A\land B]^1}{A}\;(\land E_1)
}{B\land A}\;(\land I).
\]
Zatem trzy kropki w pierwszym zapisie zastępowały
trzy konkretne kroki: uzyskanie $B$, uzyskanie $A$
i ich połączenie w $B\land A$.
\emph{W żadnym miejscu nie założyliśmy dodatkowo $B\land A$}.
Aby otrzymać dowód bez otwartych założeń,
można dopiero teraz rozładować $[A\land B]^1$
regułą $(\to I)_1$, otrzymując $(A\land B)\to(B\land A)$.
\end{przyklad}

\begin{uwaga}[Kropki nie zmieniają statusu założenia]
Na każdym etapie pominiętego fragmentu trzeba odróżnić
\emph{wnioski wyprowadzone} od \emph{założeń nadal otwartych}.
Jeżeli we fragmencie powyższego przykładu wykorzystaliśmy
$[A\land B]^1$, to otrzymane $B\land A$ wciąż zależy
od tego założenia. Sam symbol $\vdots$ go nie rozładowuje.
Rozładowanie następuje dopiero przez regułę z wyraźnym
oznaczeniem indeksu, taką jak $(\to I)_1$.
\end{uwaga}

\begin{przyklad}[Czego kropki \emph{nie} usprawiedliwiają]
Gdybyśmy bez dodatkowych założeń napisali
\[
\begin{array}{c}
A\\
\vdots\\
B
\end{array}
\]
nie wynikałoby z tego, że $A\vdash_{\mathrm{ND}}B$.
Na przykład z samego „$2$ jest liczbą parzystą”
nie da się w logice wyprowadzić dowolnego zdania
„$3$ jest liczbą parzystą”. Aby uznać taki zapis
za rzeczywisty dowód, trzeba podać pominięte kroki
lub odesłać do \emph{wcześniej udowodnionego} twierdzenia.
W szczególności $\vdots$ nie jest odpowiednikiem
przyzwolenia „tu dzieje się coś oczywistego”.
\end{przyklad}

\begin{przyklad}[Kropki w dwóch niezależnych gałęziach]
W regule eliminacji alternatywy pojawiają się dwa
poddowody, na ogół o \emph{różnej} liczbie kroków:
\[
\frac{A\lor B\qquad
  \begin{array}{c}[A]^i\\\vdots\\ C\end{array}
  \qquad
  \begin{array}{c}[B]^j\\\vdots\\ C\end{array}
}{C}\;(\lor E)_{i,j}.
\]
Pierwsza para kropek zastępuje dowód $C$ przy założeniu $A$,
a druga --- dowód tego samego $C$ przy założeniu $B$.
Nie zakładamy, że oba fragmenty są identyczne.
Na końcu rozładowujemy odpowiednio $A$ i $B$;
pozostałe założenia z obu gałęzi nadal pozostają otwarte.
Warunkiem stosowania reguły jest istnienie \emph{obu}
fragmentów, nie tylko jednej gałęzi.
Dlaczego? Z $A\lor B$ nie wiemy, który przypadek zachodzi.
Jeżeli tylko z $A$ udowodniliśmy $C$, rzeczywisty przypadek
może być $B$ i wówczas $C$ nie musi zachodzić.
Przykładowo przy wartościach $A=1$, $B=0$, $C=0$
alternatywa $A\lor B$ jest prawdziwa, ale z samego
warunkowego rozumowania „gdyby $B$, to $C$”
nie można otrzymać $C$. Oba poddowody są potrzebne.
Natomiast nie zakładamy, że \emph{równocześnie} zachodzą $A$ i $B$:
to są dwa oddzielne, warunkowe rozumowania.
\end{przyklad}

\begin{uwaga}[Trzy różne rodzaje wielokropków]
Warto rozróżniać $\vdots$ (pionowe kropki, oznaczające
pominięty fragment \emph{dedukcji}), $\cdots$ (poziomy
wielokropek, np. między kilkoma przesłankami lub elementami
ciągu) oraz $\ddots$ (kropki ukośne, używane np.
w macierzach). To tylko konwencje \emph{rysowania};
żaden z tych symboli nie zastępuje logicznej reguły.
\end{uwaga}

\section{Założenia, wnioski i rozładowywanie założeń}
\begin{definicja}[Dedukcja i dowodliwość]\label{def:nd-derivation}
Ustalmy język pierwszego rzędu $\mathcal L$.
\emph{Dedukcją} formuły $\varphi$ z założeń $\Gamma$ nazywamy
skończone drzewo formuł, którego liście są dozwolonymi
założeniami (z $\Gamma$ albo tymczasowymi) lub instancjami
reguł bez przesłanek, a każda pozostała formuła wynika
z formuł bezpośrednio ponad nią na mocy jednej z reguł
podanych poniżej. Tymczasowe założenia mogą być
\emph{rozładowane} przez odpowiednią regułę: od tej chwili
nie należą do otwartych założeń dedukcji.

Zapis $\Gamma\vdash_{\mathrm{ND}}\varphi$ oznacza, że istnieje
taka dedukcja z wnioskiem $\varphi$, której wszystkie
\emph{nierozładowane} założenia należą do $\Gamma$.
Dla $\Gamma=\varnothing$ piszemy $\vdash_{\mathrm{ND}}\varphi$.
\end{definicja}

\begin{przyklad}[Co oznacza rozładowanie założenia?]\label{ex:nd-discharge-explained}
Załóżmy, że chcemy udowodnić implikację „jeśli $A$, to $A$”.
Wewnątrz dowodu możemy na próbę założyć $A$. Skoro założyliśmy $A$,
mamy już $A$ jako wniosek tego krótkiego poddowodu. Teraz kończymy
okres obowiązywania założenia i zapisujemy implikację $A\to A$:
\[
\frac{[A]^1}{A\to A}\;(\to I)_1.
\]
Liczba $1$ przy nawiasie i przy nazwie reguły wiąże te dwa miejsca.
Nie znaczy ona, że formuła $A$ została udowodniona bez założeń.
Bez założeń otrzymaliśmy dopiero \emph{implikację} $A\to A$.
W dłuższym dowodzie wewnątrz rozładowywanego fragmentu
mogą znajdować się inne kroki i kilka użyć tego samego założenia.
\end{przyklad}


\subsection{Co oznacza nawias przy założeniu?}\label{sec:nd-assumption-scope}
\begin{uwaga}[Założenie nie jest stwierdzeniem prawdziwości]\label{rem:nd-assumption-not-truth}
Napis $[A]^i$ czytamy: „\emph{na potrzeby tego poddowodu zakładamy $A$;
założenie oznaczamy numerem $i$}”. Nawias kwadratowy nie jest spójnikiem
logicznym, a $i$ nie jest wykładnikiem potęgi. To tylko oznaczenie
\emph{zależności dowodu od przyjętego założenia}.

Nie twierdzimy, że $A$ jest prawdziwe „wewnątrz nawiasu” ani że
staje się fałszywe po wyjściu z nawiasu. W ogóle nie ustalamy tutaj
wartości logicznej $A$. W poddowodzie \emph{wolno posługiwać się $A$}
bez wcześniejszego uzasadnienia; po jego rozładowaniu nie wolno
powoływać się na \emph{to konkretne tymczasowe założenie}.
Jeżeli jednak $A$ zostało niezależnie dowiedzione albo pozostaje
jednym z trwałych założeń $\Gamma$, to możemy nadal z niego korzystać.
\end{uwaga}

\begin{przyklad}[Założenie może być nawet fałszywe]\label{ex:nd-false-assumption}
Niech $A$ oznacza zdanie „liczba $2$ jest nieparzysta”.
Wiemy, że $A$ jest fałszywe, ale nadal wolno rozważyć
\emph{hipotetyczny} dowód rozpoczynający się od $[A]^3$.
Z tego założenia otrzymujemy samo $A$; rozładowanie daje
\[
\frac{[A]^3}{A\to A}\;(\to I)_3.
\]
Nie udowodniliśmy w ten sposób, że $2$ jest nieparzyste.
Udowodniliśmy jedynie zdanie „\emph{jeśli $2$ jest nieparzyste,
to $2$ jest nieparzyste}”, czyli implikację $A\to A$.
Jest ona prawdziwa również wtedy, gdy $A$ jest fałszywe.
Właśnie dlatego rozładowanie założenia nie może oznaczać
zmiany jego wartości logicznej.
\end{przyklad}

\begin{przyklad}[Jak zmienia się zbiór otwartych założeń]\label{ex:nd-open-assumptions}
Ustalmy dowolne formuły $A$ i $B$.
Zademonstrujemy kolejno dwa zagnieżdżone poddowody:
\[
\frac{
 \displaystyle\frac{
   \displaystyle\frac{[A]^1\qquad[B]^2}{A\land B}\;(\land I)
 }{B\to(A\land B)}\;(\to I)_2
}{A\to\bigl(B\to(A\land B)\bigr)}\;(\to I)_1.
\]
Odczytujemy drzewo \emph{od góry do dołu}.
Na początku założyliśmy jednocześnie $A$ (etykieta $1$)
i $B$ (etykieta $2$), dlatego wolno nam utworzyć $A\land B$.
W tym momencie wniosek zależy od \emph{obu} tymczasowych założeń.

Teraz stosujemy $(\to I)_2$. Zamykamy założenie oznaczone $2$,
więc otrzymujemy $B\to(A\land B)$, ale nadal korzystamy z $[A]^1$:
\emph{otwarte pozostaje tylko założenie $A$}.
Dopiero zastosowanie $(\to I)_1$ zamyka także założenie $A$.
Końcowa formuła $A\to(B\to(A\land B))$ nie zależy już
od żadnego otwartego założenia, czyli
$\vdash_{\mathrm{ND}}A\to(B\to(A\land B))$.

Nie otrzymujemy natomiast ani $\vdash_{\mathrm{ND}}A$,
ani $\vdash_{\mathrm{ND}}B$. Po zakończeniu poddowodu
nie wolno po prostu „wyciągnąć” $A$ albo $B$ jako wniosku
niezależnego od założeń.
\end{przyklad}

\begin{uwaga}[Rozładowanie przy dowodzie nie wprost]
Podobnie działa reguła wprowadzania negacji:
\[
\frac{[A]^i\quad\vdots\quad\bot}{\neg A}\;(\neg I)_i.
\]
Jeżeli przyjęcie $A$ prowadzi do sprzeczności, to po
zamknięciu tego założenia możemy wnioskować $\neg A$.
To kolejny argument, że zapis $[A]^i$ \emph{nie ogłasza
prawdziwości $A$}: wewnątrz hipotetycznego poddowodu
możemy dojść właśnie do wniosku, że $A$ jest niemożliwe
(przy pozostałych otwartych założeniach).
\end{uwaga}

Przejście od poddowodu do implikacji streszcza zatem schemat
\[
\Gamma\cup\{A\}\vdash_{\mathrm{ND}}B
\quad\Longrightarrow\quad
\Gamma\vdash_{\mathrm{ND}}A\to B,
\]
gdzie $\Gamma$ oznacza zbiór pozostałych, dostępnych założeń.
Jego pełne uzasadnienie podamy w
Twierdzeniu~\ref{thm:nd-deduction}.

Ściśle rzecz biorąc, w drzewie trzeba odnotować, które
wystąpienia założeń zostają rozładowane. Będziemy używać indeksów
$[\varphi]^i$ oraz zapisu „rozładowujemy $i$”.

Dla wygody symbol $\bot$ będzie oznaczał ustaloną zamkniętą
formułę fałszywą w każdej strukturze, na przykład
$\neg\forall v_0\,(v_0=v_0)$.
Nie dodajemy więc nowego symbolu do języka: $\bot$ jest skrótem.
Znaczenie tego skrótu jest uzasadnione
Definicją~\ref{def:sem-satisfaction} i zwykłą własnością równości.

\section{Reguły dla koniunkcji i implikacji}
\begin{definicja}[Reguły $\land$ oraz $\to$]\label{def:nd-rules-conj-imp}
W dedukcji naturalnej dopuszczamy następujące schematy:
\[
\frac{\varphi\qquad\psi}{\varphi\land\psi}\;(\land I),
\qquad
\frac{\varphi\land\psi}{\varphi}\;(\land E_1),
\qquad
\frac{\varphi\land\psi}{\psi}\;(\land E_2),
\]
\[
\frac{[\varphi]^i\quad\vdots\quad\psi}
     {\varphi\to\psi}\;(\to I)_i,
\qquad
\frac{\varphi\to\psi\qquad\varphi}{\psi}\;(\to E).
\]
Reguła $(\to I)_i$ rozładowuje oznaczone wystąpienia
założenia $\varphi$ użyte w poddedukcji $\psi$.
\end{definicja}
\begin{przyklad}[Pierwszy formalny dowód]\label{ex:nd-conj-swap}
Udowodnimy $\vdash_{\mathrm{ND}}(P\land Q)\to(Q\land P)$.
Rozpoczynamy od tymczasowego założenia $P\land Q$.
Reguła $\land E_2$ daje $Q$, a reguła $\land E_1$ daje $P$.
Reguła $\land I$ pozwala połączyć te dwa wnioski w $Q\land P$.
Następnie $\to I$ rozładowuje początkowe założenie,
pozostawiając bezzałożeniowy dowód implikacji.
Pełne drzewo (a nie tylko jego fragment) ma postać
\[
\frac{\displaystyle
 \frac{
 \frac{[P\land Q]^1}{Q}\;(\land E_2)
 \qquad
 \frac{[P\land Q]^1}{P}\;(\land E_1)
 }{Q\land P}\;(\land I)
}{(P\land Q)\to(Q\land P)}\;(\to I)_1.
\]
Sprawdzamy je od góry: dwa górne liście mają tę samą postać $P\land Q$.
Z lewego odczytujemy $Q$, z prawego $P$. Łączymy je w $Q\land P$.
Na końcu zamykamy tymczasowe założenie o numerze $1$,
co pozwala wprowadzić implikację. Pod kreską końcową nie ma
już żadnego nierozładowanego założenia, stąd znak $\vdash_{\mathrm{ND}}$.
\end{przyklad}

\begin{twierdzenie}[Twierdzenie o dedukcji dla dedukcji naturalnej]
\label{thm:nd-deduction}
Dla dowolnego zbioru formuł $\Gamma$ i dowolnych formuł
$\varphi,\psi$ zachodzi
\[
\Gamma\cup\{\varphi\}\vdash_{\mathrm{ND}}\psi
\quad\Longleftrightarrow\quad
\Gamma\vdash_{\mathrm{ND}}\varphi\to\psi.
\]
\end{twierdzenie}
\begin{proof}
Załóżmy najpierw, że istnieje dedukcja $\psi$,
której wszystkie otwarte założenia należą do
$\Gamma\cup\{\varphi\}$. Jeśli niektóre otwarte liście
są wystąpieniami $\varphi$, oznaczamy je wspólnym indeksem
$i$. Zastosowanie reguły $\to I_i$ do całej dedukcji
rozładowuje te wystąpienia i daje $\varphi\to\psi$.
Pozostałe otwarte założenia pochodzą z $\Gamma$.
Nawet jeśli w poddedukcji nie było otwartego wystąpienia
$\varphi$, reguła $\to I$ jest dozwolona:
implikacja z dodatkowym, nieużytym założeniem też ma dowód.

W przeciwną stronę załóżmy, że
$\Gamma\vdash_{\mathrm{ND}}\varphi\to\psi$.
Do otrzymanego drzewa dopisujemy dodatkowy liść
$\varphi$ i stosujemy $\to E$. Wniosek to $\psi$,
a otwarte założenia należą do $\Gamma\cup\{\varphi\}$.
\end{proof}

\section{Alternatywa, równoważność i negacja}
\begin{definicja}[Pozostałe reguły zdaniowe]\label{def:nd-rules-prop}
Dla alternatywy przyjmujemy
\[
\frac{\varphi}{\varphi\lor\psi}\;(\lor I_1),
\quad
\frac{\psi}{\varphi\lor\psi}\;(\lor I_2),
\quad
\frac{\varphi\lor\psi\qquad
[\varphi]^i\;\vdots\;\theta\qquad
[\psi]^j\;\vdots\;\theta}{\theta}\;(\lor E)_{i,j}.
\]
Ostatnia reguła rozładowuje założenie $\varphi$
w pierwszej poddedukcji oraz założenie $\psi$
w drugiej. Jej sens jest następujący: rozpatrujemy oba
możliwe przypadki i dochodzimy do tego samego wniosku.

Dla negacji i sprzeczności przyjmujemy
\[
\frac{\varphi\qquad\neg\varphi}{\bot}\;(\neg E),
\quad
\frac{[\varphi]^i\;\vdots\;\bot}{\neg\varphi}\;(\neg I)_i,
\quad
\frac{\bot}{\theta}\;(\bot E).
\]
Ponieważ używamy logiki klasycznej, dodajemy regułę
\emph{dowodu nie wprost}:
\[
\frac{[\neg\varphi]^i\;\vdots\;\bot}{\varphi}\;(\mathrm{RAA})_i.
\]
Dla równoważności reguły zapisujemy jawnie:
\[
\frac{\varphi\to\psi\qquad\psi\to\varphi}
     {\varphi\leftrightarrow\psi}\;(\leftrightarrow I),
\quad
\frac{\varphi\leftrightarrow\psi}{\varphi\to\psi}\;(\leftrightarrow E_1),
\quad
\frac{\varphi\leftrightarrow\psi}{\psi\to\varphi}\;(\leftrightarrow E_2).
\]
Pierwsza reguła wymaga obu implikacji, druga i trzecia wydobywają
z równoważności odpowiednio pierwszy lub drugi kierunek.
\end{definicja}
\begin{uwaga}[Klasyczność ma znaczenie]
Reguła $\mathrm{RAA}$ nie jest regułą intuicjonistycznej
dedukcji naturalnej. Odłożenie jej prowadzi do słabszego
pojęcia wyprowadzalności. Nie będziemy zamieniać tych
dwóch rachunków bez wyraźnego zaznaczenia zmiany założeń.
\end{uwaga}

\begin{twierdzenie}[Prawo wyłączonego środka]
\label{thm:nd-lem}
Dla każdej formuły $\varphi$ zachodzi
$\vdash_{\mathrm{ND}}\varphi\lor\neg\varphi$.
\end{twierdzenie}
\begin{proof}
Załóżmy tymczasowo $\neg(\varphi\lor\neg\varphi)$.
Jeśli dodatkowo założymy $\varphi$, to przez $\lor I_1$
otrzymamy $\varphi\lor\neg\varphi$.
Razem z pierwszym założeniem daje to $\bot$ przez $\neg E$.
Rozładowując założenie $\varphi$ regułą $\neg I$,
otrzymujemy $\neg\varphi$.
Stosując teraz $\lor I_2$, otrzymujemy ponownie
$\varphi\lor\neg\varphi$, sprzeczne z pierwszym założeniem.
Zatem z $\neg(\varphi\lor\neg\varphi)$ wyprowadziliśmy $\bot$.
Reguła $\mathrm{RAA}$ rozładowuje to założenie
i daje $\varphi\lor\neg\varphi$ bez żadnych otwartych założeń.
\end{proof}

\section{Kwantyfikatory i warunki na zmienne}
Na tym etapie ujawnia się ważna różnica między rachunkiem zdań
a logiką pierwszego rzędu. Nie każda zmienna może być
uogólniana, a podstawienie nie może przechwytywać zmiennych.
Zlekceważenie tych warunków prowadziłoby do fałszywych twierdzeń.

\begin{definicja}[Reguły kwantyfikatorów]\label{def:nd-rules-quant}
Niech $t$ będzie termem dopuszczalnym do podstawienia
za $x$ w $\varphi$ (Definicja~\ref{def:syntax-free-for}).
Stosujemy reguły
\[
\frac{\forall x\,\varphi}{\varphi[x:=t]}\;(\forall E),
\qquad
\frac{\varphi}{\forall x\,\varphi}\;(\forall I),
\]
\[
\frac{\varphi[x:=t]}{\exists x\,\varphi}\;(\exists I).
\]
Reguła $\forall I$ jest dozwolona tylko wtedy, gdy $x$
nie występuje wolno w \emph{żadnym otwartym założeniu}
poddedukcji kończącej się na $\varphi$.

W regule eliminacji kwantyfikatora egzystencjalnego
stosujemy zmienną $y$, niewystępującą w ogóle
w $\varphi$ ani w końcowym wniosku $\theta$,
a ponadto niewystępującą wolno w żadnym innym
otwartym założeniu rzeczywiście użytym w obu
gałęziach dowodu (poza rozładowywanym założeniem
$[\varphi[x:=y]]^i$):
\[
\frac{\exists x\,\varphi\qquad
[\varphi[x:=y]]^i\;\vdots\;\theta}
{\theta}\;(\exists E)_i.
\]
Reguła rozładowuje tymczasowe założenie
$\varphi[x:=y]$. Warunek świeżości dotyczy obu gałęzi:
$y$ nie może występować wolno ani w pozostałych otwartych
założeniach poddedukcji, ani w otwartych założeniach dowodu
przesłanki $\exists x\,\varphi$.
Nie wymagamy natomiast, aby $y$ nie występowała w żadnej
formule całego zbioru $\Gamma$: zbiór ten może być nieskończony,
a konkretna dedukcja korzysta tylko ze skończenie wielu
otwartych założeń. W przeciwnym razie połączenie obu gałęzi
mogłoby przenieść informację o konkretnym $y$ do wniosku,
chociaż świadek kwantyfikatora egzystencjalnego nie musi
być wartością $y$. Przed zastosowaniem reguły w razie
potrzeby zmieniamy nazwy zmiennych związanych,
aby podstawienie było dopuszczalne.
\end{definicja}

\begin{przyklad}[Jak czytać reguły kwantyfikatorów]
Jeżeli znamy $\forall x\,P(x)$, wolno wstawić np. term $0$ za $x$
i otrzymać $P(0)$; taki krok nosi nazwę $(\forall E)$.
Jeżeli natomiast ustaliliśmy $P(0)$, otrzymujemy przez $(\exists I)$
formułę $\exists x\,P(x)$: liczba oznaczona przez term $0$ jest
konkretnym świadkiem istnienia. Te kroki pokazuje zapis
\[
\frac{\forall x\,P(x)}{P(0)}\;(\forall E),
\qquad
\frac{P(0)}{\exists x\,P(x)}\;(\exists I).
\]
Nie wolno jednak odwrócić pierwszej kreski: z jednego przypadku
$P(0)$ nie wynika twierdzenie o wszystkich elementach.
\end{przyklad}

\begin{przyklad}[Co znaczą kropki w eliminacji egzystencjalnej]
Załóżmy, że wiadomo tylko $\exists x\,P(x)$ i chcemy pokazać
zdanie $\theta$ niezawierające zmiennej $y$.
Tymczasowo przyjmujemy $P(y)$, gdzie $y$ jest nową nazwą,
nieobecną w innych otwartych założeniach i we wniosku $\theta$.
Jeżeli z tego założenia potrafimy wykonać jakieś kroki i dojść
do $\theta$, w schemacie oznaczamy je pionowymi kropkami:
\[
\frac{\exists x\,P(x)\qquad [P(y)]^2\;\vdots\;\theta}
{\theta}\;(\exists E)_2.
\]
Kropki nie pozwalają dowolnie przeskakiwać kroków: w rzeczywistym dowodzie
muszą odpowiadać konkretnemu, sprawdzalnemu drzewu
(zob. sekcję~\ref{sec:nd-vertical-dots}). Indeks $2$ wskazuje,
że założenie $P(y)$ po użyciu reguły nie pozostaje otwarte.
\end{przyklad}

\begin{przyklad}[Dlaczego ograniczenie dla $\forall I$ jest konieczne]
Z założenia $P(x)$ nie wolno wywnioskować $\forall x\,P(x)$.
W strukturze o uniwersum $\{0,1\}$ niech $P$ oznacza zbiór $\{0\}$.
Przy wartościowaniu $s(x)=0$ założenie $P(x)$ jest prawdziwe,
a $\forall x\,P(x)$ jest fałszywe. Niepoprawny krok
polegałby na zastosowaniu $\forall I$ do zmiennej wolnej
w otwartym założeniu $P(x)$.
\end{przyklad}

\begin{przyklad}[Dlaczego świadek egzystencjalny jest lokalny]
Z $\exists x\,P(x)$ nie możemy po prostu wyprowadzić
$P(y)$, jeśli $y$ jest ustaloną wcześniej zmienną.
W strukturze o dwóch elementach, przy $P$ prawdziwym
tylko dla jednego z nich, poprzednik jest prawdziwy,
a $P(y)$ może być fałszywe dla wartościowania przypisującego
$y$ drugi element. Reguła $\exists E$ pozwala użyć
świeżej zmiennej \emph{wewnątrz poddedukcji}, ale nie
pozwala jej przemycić do końcowego wniosku.
\end{przyklad}

\begin{twierdzenie}[Przykład dowodzenia z kwantyfikatorami]
\label{thm:nd-univ-chain}
Dla formuł jednoargumentowych $P(x)$ i $Q(x)$ zachodzi
\[
\bigl\{\forall x\,(P(x)\to Q(x)),\ \forall x\,P(x)\bigr\}
\vdash_{\mathrm{ND}}\forall x\,Q(x).
\]
\end{twierdzenie}
\begin{proof}
Oznaczmy oba założenia odpowiednio przez $H_1$ i $H_2$.
Zastosujmy $\forall E$, podstawiając za $x$ samą
zmienną $x$. Takie podstawienie jest zawsze dopuszczalne,
nawet gdy $P(x)$ lub $Q(x)$ zawierają wewnętrzne
kwantyfikatory wiążące inne wystąpienia zmiennych.
Z $H_1$ otrzymujemy $P(x)\to Q(x)$,
a z $H_2$ otrzymujemy $P(x)$.
Reguła $\to E$ daje $Q(x)$.
W żadnym z otwartych założeń $H_1,H_2$ zmienna $x$
nie występuje wolno: wszystkie jej wystąpienia
w tych zdaniach są związane przez zewnętrzne kwantyfikatory.
Możemy zatem zastosować $\forall I$ i otrzymać
bezpośrednio $\forall x\,Q(x)$.
\end{proof}

\begin{twierdzenie}[Od istnienia przypadku do istnienia własności]
\label{thm:nd-exists-projection}
Niech $P(x)$ i $Q(x)$ będą formułami o jedynej zmiennej wolnej $x$.
Wtedy
\[
\exists x\,(P(x)\land Q(x))
\vdash_{\mathrm{ND}}\exists x\,P(x).
\]
\end{twierdzenie}
\begin{proof}
Przyjmijmy założenie $H:=\exists x\,(P(x)\land Q(x))$.
Wybierzmy nową zmienną $y$, niewystępującą w $H$ ani
we wniosku. W poddedukcji przyjmijmy tymczasowo
$P(y)\land Q(y)$. Reguła $\land E_1$ daje $P(y)$,
a następnie $\exists I$ daje $\exists x\,P(x)$.
Zastosowanie $\exists E$ rozładowuje założenie
$P(y)\land Q(y)$; wniosek nie zawiera wolnego $y$,
więc jest dozwolony. Jedynym pozostałym założeniem jest $H$.
\end{proof}

\begin{przyklad}[Dlaczego świeżość obejmuje także przesłankę główną]
\label{ex:nd-exists-major-premise-fresh-46}
Najpierw rozważmy błędne użycie zmiennej występującej wolno
w \emph{innym otwartym założeniu}. Niech $P$ będzie
predykatem jednoargumentowym, a $R$ skrótem zdania
$\forall z\,P(z)$. Z założeń
$\exists x\,P(x)$ oraz $P(y)\to R$ nie wolno
wywnioskować $R$ przez poddedukcję od $[P(y)]^i$:
chociaż z $P(y)$ i drugiego założenia otrzymujemy $R$,
zmienna $y$ pozostaje wolna w $P(y)\to R$.

W uniwersum $\{0,1\}$ z $P^{\mathfrak A}=\{0\}$
oraz $s(y)=1$ obydwa założenia są spełnione,
lecz $R=\forall z\,P(z)$ jest fałszywe.
To kontrmodel dla takiego niepoprawnego wnioskowania.

Osobno sprawdźmy konieczność świeżości względem
\emph{przesłanki głównej}. Niech teraz
$H:=\exists x\,\neg(x=y)$, a jako żądany wniosek
weźmy fałszywe zdanie
$\theta:=\forall z\,\neg(z=z)$.
Gdyby wolno było użyć $y$ jako zmiennej świadka,
otrzymalibyśmy tymczasowe założenie
$[\neg(y=y)]^i$. Reguła $=I$ daje $y=y$,
więc $\neg E$ daje $\bot$, a następnie $\bot E$
pozwala wyprowadzić $\theta$.
Taki schemat jest \emph{niedozwolony}, ponieważ $y$
występuje wolno w $H$.
W dwuelementowym uniwersum $H$ jest prawdziwe
przy każdym wartościowaniu $y$ (istnieje inny element),
lecz $\theta$ jest fałszywe.
Oba kontrprzykłady pokazują, że świeżość
jest konieczna do zachowania prawdy.
\end{przyklad}

\section{Równość i podstawianie równych obiektów}
\begin{definicja}[Reguły równości]\label{def:nd-equality}
Dopuszczamy regułę $=I$: dla dowolnego termu $t$
można bez założeń wyprowadzić $t=t$. Zapisujemy ją jako
$\frac{}{t=t}\;(=I)$. Puste miejsce nad kreską znaczy,
że ta reguła nie wymaga żadnej wcześniejszej przesłanki.
Reguła $=E$ ma postać: z $t=u$ i
$\theta[x:=t]$ wolno wyprowadzić $\theta[x:=u]$,
o ile obydwa podstawienia są dopuszczalne.
\end{definicja}
\begin{przyklad}[Symetria równości]
Z założenia $t=u$ chcemy wyprowadzić $u=t$.
Najpierw przez $=I$ otrzymujemy $t=t$.
Weźmy formułę $\theta(x)$ postaci $x=t$.
Jej instancją $\theta[x:=t]$ jest $t=t$,
a instancją $\theta[x:=u]$ jest $u=t$.
Z przesłanek $t=u$ i $t=t$ reguła $=E$ daje $u=t$.
\end{przyklad}
\begin{przyklad}[Niesprzeczność wynika z prawdziwości równości]
Dla każdego termu $t$ mamy bezzałożeniowy dowód $t=t$.
Gdyby z jakichś założeń uzyskano $\neg(t=t)$,
reguła $\neg E$ po połączeniu z $t=t$ dawałaby $\bot$.
Jest to prosty przykład konfliktu formalnego z podstawowym
prawem logiki z równością.
\end{przyklad}

\section{Własność skończonego użycia założeń}
\begin{twierdzenie}[Skończona liczba użytych przesłanek]
\label{thm:nd-finite-support}
Jeżeli $\Gamma\vdash_{\mathrm{ND}}\varphi$, to istnieje
skończony zbiór $\Gamma_0\subseteq\Gamma$ taki, że
$\Gamma_0\vdash_{\mathrm{ND}}\varphi$.
\end{twierdzenie}
\begin{proof}
Dowód formalny jest z definicji skończonym drzewem.
Ma więc tylko skończenie wiele liści. Wybierzmy spośród
nich te, które pozostają otwartymi założeniami, i zbierzmy
występujące w nich formuły do zbioru $\Gamma_0$.
Zbiór ten jest skończony i zawiera się w $\Gamma$.
To samo drzewo, z tymi samymi regułami i rozładowaniami,
jest dedukcją $\varphi$ z $\Gamma_0$.
\end{proof}

\begin{uwaga}[Czego jeszcze nie dowiedliśmy?]
Sam fakt, że zapisaliśmy reguły, nie rozstrzyga,
czy zawsze prowadzą one od zdań prawdziwych do prawdziwych.
Tę własność — \emph{poprawność} rachunku dowodowego —
udowodnimy w następnym rozdziale.
Nie udowodniliśmy także, że z każdej konsekwencji
semantycznej można otrzymać formalny dowód.
Taką tezę nazywamy \emph{pełnością} i będzie ona przedmiotem
osobnego, znacznie trudniejszego argumentu.
\end{uwaga}